% =============================================================================
% ESTADÍSTICA CON R: FUNDAMENTOS Y APLICACIONES
% Libro técnico — clase book de LaTeX
% Compilar con: pdflatex main.tex (2 pasadas) o usar compile.sh
% =============================================================================

\documentclass[
  10pt,           % Tamaño de fuente base (estándar académico compacto)
  twoside,        % Impresión doble cara
  openright       % Capítulos empiezan en página derecha
]{book}

% Incluir preámbulo (paquetes, estilos, comandos)
% =============================================================================
% PREÁMBULO: Estadística con R — Fundamentos y Aplicaciones
% =============================================================================

% --- Codificación y lenguaje (LuaLaTeX maneja UTF-8 nativamente) -------------
\usepackage{fontspec}         % Gestión de fuentes OpenType con LuaLaTeX
\defaultfontfeatures{Ligatures=TeX}

% Babel v3.x (TeX Live 2025): carga idioma con \babelprovide
\usepackage{babel}
\babelprovide[main, import]{spanish}

% --- Tipografía --------------------------------------------------------------
\usepackage{lmodern}          % Latin Modern
\usepackage{microtype}        % Optimización tipográfica (kerning, protrusion)

% --- Matemáticas -------------------------------------------------------------
\usepackage{amsmath}          % Entornos equation, align, gather, etc.
\usepackage{amssymb}          % Símbolos matemáticos adicionales
\usepackage{amsthm}           % Entornos theorem, proof, definition
\usepackage{bm}               % Negritas en modo matemático (\bm)
\usepackage{mathtools}        % Extensiones de amsmath

% --- Definiciones de entornos matemáticos -----------------------------------
\theoremstyle{definition}
\newtheorem{definicion}{Definición}[chapter]
\newtheorem{ejemplo}{Ejemplo}[chapter]
\newtheorem{propiedad}{Propiedad}[chapter]

\theoremstyle{plain}
\newtheorem{teorema}{Teorema}[chapter]
\newtheorem{corolario}{Corolario}[chapter]
\newtheorem{lema}{Lema}[chapter]

\theoremstyle{remark}
\newtheorem{nota}{Nota}[chapter]
\newtheorem{observacion}{Observación}[chapter]

% --- Gráficos y figuras ------------------------------------------------------
\usepackage{graphicx}         % Inclusión de imágenes
\usepackage{float}            % Control de posición [H]
\usepackage{subcaption}       % Subfiguras
\graphicspath{{../output/}{./figuras/}}

% --- Tablas ------------------------------------------------------------------
\usepackage{booktabs}         % Líneas profesionales (\toprule, \midrule, \bottomrule)
\usepackage{array}            % Extensiones de columnas
\usepackage{longtable}        % Tablas que cruzan páginas (generadas por Pandoc)
\usepackage{multirow}         % Celdas que abarcan filas

% --- Código fuente R ---------------------------------------------------------
\usepackage{listings}
\usepackage{xcolor}

% Colores para el resaltado de sintaxis
\definecolor{codebg}{RGB}{248, 248, 248}
\definecolor{codeframe}{RGB}{200, 200, 200}
\definecolor{codecomment}{RGB}{106, 153, 85}
\definecolor{codekeyword}{RGB}{86, 156, 214}
\definecolor{codestring}{RGB}{206, 145, 120}
\definecolor{codenumber}{RGB}{181, 206, 168}
\definecolor{rblue}{RGB}{22, 92, 170}

% Estilo de código R
\lstdefinestyle{estilo_r}{
  language        = R,
  basicstyle      = \ttfamily\scriptsize,
  backgroundcolor = \color{codebg},
  frame           = single,
  framerule       = 0.5pt,
  rulecolor       = \color{codeframe},
  commentstyle    = \color{codecomment}\itshape,
  keywordstyle    = \color{codekeyword}\bfseries,
  stringstyle     = \color{codestring},
  numberstyle     = \tiny\color{codenumber},
  numbers         = left,
  numbersep       = 8pt,
  stepnumber      = 1,
  tabsize         = 2,
  showstringspaces= false,
  breaklines      = true,
  breakatwhitespace= true,
  captionpos      = b,
  morekeywords    = {TRUE, FALSE, NULL, NA, NaN, Inf,
                     library, require, install.packages,
                     function, return, if, else, for, while, repeat,
                     in, next, break,
                     data.frame, list, vector, matrix, array, factor,
                     read.csv, write.csv, source, setwd, getwd,
                     ggplot, aes, geom_point, geom_line, geom_bar,
                     geom_histogram, geom_boxplot, geom_smooth,
                     labs, theme_minimal, ggsave, scale_color_viridis_d,
                     lm, glm, aov, summary, coef, residuals, fitted,
                     t.test, wilcox.test, shapiro.test, cor, var.test,
                     mean, median, sd, var, IQR, quantile, range,
                     table, prop.table, tapply, sapply, lapply, apply,
                     filter, select, mutate, group_by, summarise, arrange,
                     pivot_wider, pivot_longer, left_join, inner_join,
                     set.seed, runif, rnorm, sample, rep, seq},
  literate        =
    {á}{{\'a}}1 {é}{{\'e}}1 {í}{{\'{\i}}}1 {ó}{{\'o}}1 {ú}{{\'u}}1
    {Á}{{\'A}}1 {É}{{\'E}}1 {Í}{{\'I}}1 {Ó}{{\'O}}1 {Ú}{{\'U}}1
    {ñ}{{\~n}}1 {Ñ}{{\~N}}1 {ü}{{\"u}}1 {Ü}{{\"U}}1
}

\lstset{style=estilo_r}

% Entorno para código inline
\newcommand{\rcode}[1]{\texttt{\small\color{rblue}#1}}

% --- Hiperenlaces y metadatos PDF -------------------------------------------
\usepackage{hyperref}
\hypersetup{
  pdftitle    = {Estadística con R: Fundamentos y Aplicaciones},
  pdfauthor   = {Proyecto Libro Caribe Colombiano},
  pdfsubject  = {Estadística, R, Análisis de Datos},
  pdfkeywords = {estadística descriptiva, inferencia, regresión, R, tidyverse,
                 Caribe colombiano, Sierra Nevada, Cesar, Barranquilla},
  colorlinks  = true,
  linkcolor   = rblue,
  urlcolor    = rblue,
  citecolor   = rblue,
  bookmarks   = true,
  bookmarksnumbered = true
}

% --- Diseño de página --------------------------------------------------------
\usepackage{geometry}
\geometry{
  letterpaper,                 % Formato carta (21.59 × 27.94 cm)
  top        = 2.5cm,
  bottom     = 2.5cm,
  left       = 3.0cm,          % Margen izquierdo / interior (encuadernación)
  right      = 2.5cm,          % Margen derecho / exterior
  headheight = 14pt
}

% --- Encabezados y pies de página -------------------------------------------
\usepackage{fancyhdr}
\pagestyle{fancy}
\fancyhf{}
\fancyhead[LE]{\small\thepage \quad \textit{Estadística con R}}
\fancyhead[RO]{\small\textit{\nouppercase{\rightmark}} \quad \thepage}
\fancyfoot[C]{}
\renewcommand{\headrulewidth}{0.4pt}

% --- Espaciado y párrafos ----------------------------------------------------
\usepackage{setspace}
\setstretch{1.15}              % Interlineado compacto (1.15x, vs 1.5x anterior)
\setlength{\parskip}{4pt}
\setlength{\parindent}{0pt}   % Sin sangría de primera línea (estilo moderno)

% --- Cajas de texto (notas, advertencias) ------------------------------------
\usepackage{tcolorbox}
\tcbuselibrary{skins, breakable}

\newtcolorbox{notabox}[1][]{
  enhanced, breakable,
  colback  = blue!5!white,
  colframe = rblue,
  title    = {#1},
  fonttitle= \bfseries\small,
  left=6pt, right=6pt, top=4pt, bottom=4pt
}

\newtcolorbox{codigobox}[1][]{
  enhanced, breakable,
  colback  = codebg,
  colframe = codeframe,
  title    = {#1},
  fonttitle= \ttfamily\small\color{rblue},
  left=4pt, right=4pt, top=2pt, bottom=2pt
}

% --- Referencias bibliográficas ----------------------------------------------
\usepackage[authoryear, round]{natbib}
\bibliographystyle{apalike}

% --- Compatibilidad con Pandoc -----------------------------------------------
% \tightlist: listas compactas generadas por Pandoc
\providecommand{\tightlist}{%
  \setlength{\itemsep}{2pt}\setlength{\parskip}{0pt}}
% \passthrough: código inline (texttt es seguro dentro de tablas)
\providecommand{\passthrough}[1]{\texttt{#1}}

% --- Comandos personalizados -------------------------------------------------
\newcommand{\E}[1]{\mathbb{E}\left[#1\right]}          % Esperanza
\newcommand{\Var}[1]{\mathrm{Var}\left(#1\right)}      % Varianza
\newcommand{\Cov}[2]{\mathrm{Cov}\left(#1,#2\right)}   % Covarianza
\newcommand{\N}[2]{\mathcal{N}\left(#2,#2\right)}      % Normal
\newcommand{\iid}{\overset{\text{iid}}{\sim}}           % iid
\newcommand{\pvalue}{p\text{-valor}}                    % p-valor
\newcommand{\dataset}[1]{\texttt{#1}}                   % Nombre de dataset


% =============================================================================
% METADATOS DEL LIBRO
% =============================================================================
\title{%
  \vspace{-2cm}
  {\Huge\bfseries Estadística Aplicada con R}\\[0.5cm]
  {\Large\itshape Fundamentos y Aplicaciones}\\[1cm]
  {\normalsize\upshape Froylan Jiménez Sánchez}\\[0.15cm]
  {\small Grupo PROPED $\cdot$ Universidad de Sucre}\\[0.5cm]
  {\normalsize\upshape Liliana Margarita Vitola Garrido}\\[0.15cm]
  {\small Universidad de Sucre}
}

\author{}

\date{%
  \small
  Primera edición $\cdot$ 2026\\[0.5cm]
  {\footnotesize Compilado con \LaTeX{} y R \texttt{4.x} $\cdot$ Pandoc 3.1}
}

% =============================================================================
% DOCUMENTO
% =============================================================================
\begin{document}

% --- Portada -----------------------------------------------------------------
\frontmatter
\maketitle
\thispagestyle{empty}

% --- Página de derechos ------------------------------------------------------
\newpage
\thispagestyle{empty}
\vspace*{\fill}
\begin{center}
  \small
  \textbf{Estadística Aplicada con R: Fundamentos y Aplicaciones}\\[0.5cm]
  Primera edición, \the\year\\[1cm]
  Este obra está bajo una licencia\\
  \textbf{Creative Commons Atribución 4.0 Internacional (CC BY 4.0)}\\[0.3cm]
  Puede compartir y adaptar el material siempre que\\
  se otorgue crédito apropiado al autor original.\\[1cm]
  {\ttfamily https://creativecommons.org/licenses/by/4.0/}\\[2cm]
  Tipografía compuesta en \LaTeX{} con la clase \texttt{book}\\
  Código R formateado con el paquete \texttt{listings}\\
  Análisis estadístico realizado con R \texttt{4.x} y Tidyverse\\
  Datasets sintéticos generados con semilla \texttt{set.seed(42/123/2024)}
\end{center}
\vspace*{\fill}

% --- Dedicatoria -------------------------------------------------------------
\newpage
\thispagestyle{empty}
\vspace*{5cm}
\begin{flushright}
  \itshape
  A los estudiantes del Caribe colombiano\\
  que convierten datos en conocimiento\\
  y conocimiento en desarrollo regional.
\end{flushright}

% --- Prefacio ----------------------------------------------------------------
\chapter*{Prefacio}
\addcontentsline{toc}{chapter}{Prefacio}

Este libro nació de una convicción simple: \textbf{la estadística se aprende mejor con datos propios}.
Por eso, cada ejemplo, cada gráfico y cada ejercicio de estas páginas está construido sobre
fenómenos reales del Caribe colombiano: la biodiversidad escalonada de la Sierra Nevada de
Santa Marta, la productividad de los cultivos de palma de aceite en el Cesar y la logística
operacional del Puerto de Barranquilla.

\medskip

\textbf{¿A quién va dirigido?} A estudiantes de pregrado y posgrado en ciencias naturales,
agronómicas, económicas e ingenierías que necesiten incorporar el análisis cuantitativo en su
práctica profesional. Se asume conocimiento de matemáticas de bachillerato, pero no de programación.

\medskip

\textbf{¿Cómo usar este libro?} Cada capítulo combina teoría con código R ejecutable.
Se recomienda tener RStudio abierto y reproducir cada bloque de código mientras se lee.
Los \emph{datasets} sintéticos del libro son completamente reproducibles: ejecutar el script
\texttt{scripts/generar\_todos\_datasets.R} recreará exactamente los mismos datos.

\medskip

\textbf{Estructura:}
\begin{itemize}
  \item \textbf{Capítulo 1} — Introduce R como entorno de trabajo estadístico.
  \item \textbf{Capítulo 2} — Estadística descriptiva: resumir y visualizar datos.
  \item \textbf{Capítulo 3} — Inferencia: generalizaciones con incertidumbre medida.
  \item \textbf{Capítulo 4} — Regresión y modelos: cuantificar relaciones entre variables.
\end{itemize}

\medskip

Todo el código fuente, los datos y este documento \LaTeX{} están disponibles en el repositorio
del proyecto bajo licencia \textbf{CC BY 4.0}.

\vspace{1cm}
\begin{flushright}
  \itshape
  Froylan Jiménez Sánchez \\ Liliana Margarita Vitola Garrido\\[0.3cm]
  \small Universidad de Sucre, 2026
\end{flushright}

% --- Tabla de contenidos -----------------------------------------------------
\tableofcontents
\listoffigures
\listoftables

% --- Sobre los Autores (páginas preliminares — numeración romana) -------------
% =============================================================================
% SOBRE LOS AUTORES
% =============================================================================

\chapter*{Sobre los Autores}
\addcontentsline{toc}{chapter}{Sobre los Autores}

\bigskip

% -----------------------------------------------------------------------------
\subsection*{Froylan Jiménez Sánchez}
% -----------------------------------------------------------------------------

\noindent
\textbf{Froylan Jiménez Sánchez} es docente de matemáticas y física con más de diez años
de experiencia en educación superior. Es Especialista en Matemática Avanzada y Magíster en
Estadística Aplicada, con investigación centrada en la intersección entre métodos
estadísticos, inteligencia artificial y sistemas agroindustriales.

\medskip

Se desempeña como docente e investigador vinculado al \textbf{Grupo PROPED}
(\textit{Grupo de Investigación en Procesos Pedagógicos y Didácticos}) de la
\textit{Universidad de Sucre} (Sincelejo, Colombia), institución donde integra el
análisis estadístico y computacional a la formación cuantitativa de estudiantes de
pregrado en ciencias naturales, ingenierías y ciencias de la salud.

\medskip

Sus áreas de investigación comprenden el aprendizaje automático (\textit{Machine Learning})
aplicado a la agricultura de precisión, la bioimpedancia espectroscópica para la evaluación
no destructiva de madurez en aguacate Hass, los \textit{Fuzzy Cognitive Maps} como
herramienta de modelado de sistemas complejos, y la estadística aplicada al análisis de
datos agronómicos y educativos.

\medskip

Entre sus contribuciones científicas más recientes se destacan un artículo publicado en
\textbf{MDPI \textit{Computers}} (2026) sobre la evaluación de la madurez del aguacate
Hass mediante bioimpedancia espectroscópica y mapas cognitivos difusos, así como un
capítulo de libro en editorial \textbf{Springer} sobre aplicaciones de Machine Learning
en sistemas agrícolas.

\bigskip\bigskip

% -----------------------------------------------------------------------------
\subsection*{Liliana Margarita Vitola Garrido}
% -----------------------------------------------------------------------------

\noindent
\textbf{Liliana Margarita Vitola Garrido} es Licenciada en Matemáticas de la
\textit{Universidad de Sucre} (2004) y Magíster en Estadística Aplicada de la
\textit{Fundación Universidad del Norte} (Barranquilla, 2012). Su trayectoria
académica combina una formación matemática rigurosa con una especialización en
métodos estadísticos orientados a la investigación aplicada.

\medskip

Actualmente se desempeña como \textbf{Docente de Estadística, Diseño de Experimentos
y Matemáticas} en la \textit{Universidad de Sucre} (Sincelejo, Colombia), institución
en la que ha contribuido a la formación cuantitativa de múltiples generaciones de
profesionales de las ciencias naturales, la salud y las ingenierías. Su labor docente
se caracteriza por vincular los fundamentos teóricos de la estadística con su
aplicación directa a problemáticas regionales concretas.

\medskip

Sus áreas de investigación incluyen la regresión logística aplicada al análisis del
rendimiento académico, la lógica difusa (\textit{Fuzzy Logic}) y las funciones de
probabilidad, la validación de instrumentos de medición en salud mental, y el
tratamiento de agua mediante coagulantes naturales como alternativa para comunidades
con acceso limitado a saneamiento básico.

\medskip

Su producción científica abarca publicaciones en revistas indexadas nacionales e
internacionales. Entre las más representativas figuran: \textit{``Fuzzy Logic and
Probability Functions''}, publicado en el \textbf{Colombia International Journal of
Mechanical Engineering and Technology} (2019); \textit{``Supply and Demand for
Drinking Water in Sincelejo City, Colombia: A Review of Alternative Solutions''},
en \textbf{IAEME Publication} (2019); \textit{``Regresión Logística: Una aplicación
en la identificación de variables que inciden en el rendimiento académico en el área
de matemáticas''}, editado por la \textbf{Universidad Militar Nueva Granada} (2015);
y \textit{``Logistic Regression: For the Identification of Socio-Economic Variables
that Influence on the Academic Performance of Students''}, en el
\textbf{Indian Journal of Science and Technology} (2017). Completa su perfil
investigativo un trabajo sobre la actividad coagulante de semillas de
\textit{Moringa oleifera} para el tratamiento de aguas residuales, que evidencia
su interés en soluciones ambientales de bajo costo para el Caribe colombiano.


% =============================================================================
% CUERPO PRINCIPAL
% =============================================================================
\mainmatter

% Restaurar estilo de encabezado para cuerpo del libro
\pagestyle{fancy}

% --- Capítulo 1 --------------------------------------------------------------
\input{capitulos/cap01}
\cleardoublepage

% --- Capítulo 2 --------------------------------------------------------------
\input{capitulos/cap02}
\cleardoublepage

% --- Capítulo 3 --------------------------------------------------------------
\input{capitulos/cap03}
\cleardoublepage

% --- Capítulo 4 --------------------------------------------------------------
\input{capitulos/cap04}
\cleardoublepage

% =============================================================================
% MATERIAL FINAL
% =============================================================================
\backmatter

% --- Apéndice A: Tablas estadísticas -----------------------------------------
\appendix
\chapter{Tablas Estadísticas de Referencia}

\section{Distribución Normal Estándar $Z$}

Valores de $P(Z \leq z)$ para la distribución normal estándar
$Z \sim \mathcal{N}(0, 1)$:

\begin{table}[H]
\centering
\caption{Percentiles de la distribución normal estándar}
\begin{tabular}{@{}lllllll@{}}
\toprule
$z$ & 0.00 & 0.01 & 0.02 & 0.03 & 0.04 & 0.05 \\
\midrule
0.0 & 0.5000 & 0.5040 & 0.5080 & 0.5120 & 0.5160 & 0.5199 \\
0.5 & 0.6915 & 0.6950 & 0.6985 & 0.7019 & 0.7054 & 0.7088 \\
1.0 & 0.8413 & 0.8438 & 0.8461 & 0.8485 & 0.8508 & 0.8531 \\
1.5 & 0.9332 & 0.9345 & 0.9357 & 0.9370 & 0.9382 & 0.9394 \\
1.6 & 0.9452 & 0.9463 & 0.9474 & 0.9484 & 0.9495 & 0.9505 \\
1.9 & 0.9713 & 0.9719 & 0.9726 & 0.9732 & 0.9738 & 0.9744 \\
2.0 & 0.9772 & 0.9778 & 0.9783 & 0.9788 & 0.9793 & 0.9798 \\
2.5 & 0.9938 & 0.9940 & 0.9941 & 0.9943 & 0.9945 & 0.9946 \\
3.0 & 0.9987 & 0.9987 & 0.9987 & 0.9988 & 0.9988 & 0.9989 \\
\bottomrule
\end{tabular}
\end{table}

\medskip
Valores críticos de uso frecuente:
\[
z_{0.90} = 1.282 \quad z_{0.95} = 1.645 \quad z_{0.975} = 1.960 \quad z_{0.995} = 2.576
\]

\section{Distribución $t$ de Student}

\begin{table}[H]
\centering
\caption{Valores críticos de la distribución $t_\nu$ (colas)}
\begin{tabular}{@{}lcccc@{}}
\toprule
$\nu$ (gl) & $t_{0.10}$ & $t_{0.05}$ & $t_{0.025}$ & $t_{0.005}$ \\
\midrule
1  & 3.078 & 6.314 & 12.706 & 63.657 \\
5  & 1.476 & 2.015 &  2.571 &  4.032 \\
10 & 1.372 & 1.812 &  2.228 &  3.169 \\
20 & 1.325 & 1.725 &  2.086 &  2.845 \\
30 & 1.310 & 1.697 &  2.042 &  2.750 \\
60 & 1.296 & 1.671 &  2.000 &  2.660 \\
$\infty$ & 1.282 & 1.645 & 1.960 & 2.576 \\
\bottomrule
\end{tabular}
\end{table}

\section{Valores críticos $\chi^2$}

\begin{table}[H]
\centering
\caption{Valores críticos de la distribución $\chi^2_\nu$}
\begin{tabular}{@{}lccccc@{}}
\toprule
$\nu$ & $\chi^2_{0.995}$ & $\chi^2_{0.975}$ & $\chi^2_{0.050}$ & $\chi^2_{0.025}$ & $\chi^2_{0.005}$ \\
\midrule
1  & 0.000 & 0.001 & 3.841 &  5.024 &  7.879 \\
3  & 0.072 & 0.216 & 7.815 &  9.348 & 12.838 \\
5  & 0.412 & 0.831 & 11.070 & 12.833 & 16.750 \\
10 & 2.156 & 3.247 & 18.307 & 20.483 & 25.188 \\
20 & 7.434 & 9.591 & 31.410 & 34.170 & 39.997 \\
\bottomrule
\end{tabular}
\end{table}

% --- Apéndice B: Comandos R esenciales ---------------------------------------
\chapter{Referencia Rápida de R}

\section{Operaciones fundamentales}

\begin{lstlisting}[caption={Comandos R de uso frecuente}]
# Ayuda
?funcion            # Documentacion de funcion
help(funcion)       # Equivalente
example(lm)         # Ejemplos ejecutables

# Entorno
ls()                # Listar objetos en memoria
rm(list = ls())     # Limpiar entorno
getwd()             # Directorio actual
setwd("ruta/")      # Cambiar directorio

# Inspeccionar objetos
str(objeto)         # Estructura
class(objeto)       # Clase/tipo
dim(objeto)         # Dimensiones [filas, columnas]
head(objeto, 10)    # Primeras 10 filas
summary(objeto)     # Resumen estadistico

# I/O de datos
df <- read.csv("archivo.csv", encoding = "UTF-8")
write.csv(df, "salida.csv", row.names = FALSE)
\end{lstlisting}

\section{Funciones estadísticas clave en R}

\begin{table}[H]
\centering
\caption{Funciones estadísticas de R ordenadas por capítulo}
\begin{tabular}{@{}lll@{}}
\toprule
Función & Descripción & Capítulo \\
\midrule
\texttt{mean(x)} & Media aritmética & 2 \\
\texttt{median(x)} & Mediana & 2 \\
\texttt{sd(x)} & Desviación estándar & 2 \\
\texttt{var(x)} & Varianza & 2 \\
\texttt{quantile(x, p)} & Cuantil de orden $p$ & 2 \\
\texttt{IQR(x)} & Rango intercuartílico & 2 \\
\texttt{table(x)} & Tabla de frecuencias & 2 \\
\texttt{cor(x, y)} & Correlación de Pearson & 4 \\
\texttt{t.test(x)} & Prueba $t$ & 3 \\
\texttt{shapiro.test(x)} & Prueba de normalidad & 3 \\
\texttt{aov(y \textasciitilde{} x)} & ANOVA de un factor & 3 \\
\texttt{lm(y \textasciitilde{} x)} & Regresión lineal & 4 \\
\texttt{glm(y \textasciitilde{} x, family)} & Modelos lineales generalizados & 4 \\
\texttt{step(modelo)} & Selección de variables & 4 \\
\bottomrule
\end{tabular}
\end{table}

% --- Bibliografía ------------------------------------------------------------
\chapter*{Referencias Bibliográficas}
\addcontentsline{toc}{chapter}{Referencias Bibliográficas}

\begin{itemize}
  \item Fedepalma (2022). \textit{Anuario Estadístico: La agroindustria de la palma de aceite en Colombia}. Federación Nacional de Cultivadores de Palma de Aceite. Bogotá.

  \item Holdridge, L. R. (1967). \textit{Life Zone Ecology}. Tropical Science Center. San José, Costa Rica.

  \item Ihaka, R., \& Gentleman, R. (1996). R: A language for data analysis and graphics. \textit{Journal of Computational and Graphical Statistics}, 5(3), 299--314.

  \item R Core Team (2024). \textit{R: A Language and Environment for Statistical Computing}. R Foundation for Statistical Computing. Vienna, Austria. \texttt{https://www.R-project.org/}

  \item Wickham, H., et al. (2019). Welcome to the Tidyverse. \textit{Journal of Open Source Software}, 4(43), 1686.

  \item Wickham, H., \& Grolemund, G. (2023). \textit{R for Data Science} (2nd ed.). O'Reilly Media. Disponible en: \texttt{https://r4ds.hadley.nz/}

  \item Montgomery, D. C., \& Runger, G. C. (2018). \textit{Applied Statistics and Probability for Engineers} (7th ed.). Wiley.

  \item Casella, G., \& Berger, R. L. (2002). \textit{Statistical Inference} (2nd ed.). Duxbury Press.

  \item James, G., Witten, D., Hastie, T., \& Tibshirani, R. (2021). \textit{An Introduction to Statistical Learning with Applications in R} (2nd ed.). Springer. Disponible en: \texttt{https://www.statlearning.com/}
\end{itemize}

% ── Apéndices ──────────────────────────────────────────────
\input{capitulos/apendice_bibliografia}
\input{capitulos/apendice_glosario}

% --- Índice analítico (requiere ejecutar makeindex) --------------------------
% \cleardoublepage
% \addcontentsline{toc}{chapter}{Índice analítico}
% \printindex

\end{document}
